\scanpage{224}
\margin{223}
\begin{center}-- 7 --\end{center}
\margin{n$^\circ$370}
si $f\le\sum_{i=1}^n c_i\gamma(s_i)g$ avec des $c_i\ge0$,
on a $\sup f\le\sum_{i=1}^n c_i g(s_i^{-1}s)$ pour un $s\in G$,
donc $\sup f\le(\sum_{i=1}^n c_i)\sup g$, d’où (iv).
Prouvons (v) ; soient $f\in\mathcal K$, $g,h\in\mathcal K_+^*$ ;
si $f\le\sum_{i=1}^n c_i\gamma(s_i)g$ et
$g\le\sum_{j=1}^p d_j\gamma(t_j)h$ ($c_i,d_j\ge0$, $s_i,t_j\in G$),
$f\le\sum_{i,j}c_i d_j\gamma(s_i t_j)h$, donc
$(f:h)\le\sum_{i,j}c_i d_j=(\sum_i c_i)(\sum_j d_j)$ ;
donc $(f:h)\le(f:g)(g:h)$.
Si on applique (v) à $f_0,f,g$ d’une part et à $f,f_0,g$ d’autre part,
on obtient (vi). Enfin, soient $f,f',h\in\mathcal K_+$ avec $h(s)\ge1$
sur le support de $f+f'$, et soit $\varepsilon>0$.
Posons $F=f+f'+\frac12\varepsilon h$ ; les fonctions $\varphi,\varphi'$,
qui coïncident respectivement avec $f/F$ et $f'/F$ sur le support de $f+f'$
et qui sont nulles en dehors de celui-ci, appartiennent à $\mathcal K_+$ ;
pour tout $\eta>0$, il existe un voisinage compact $V$ de $e$ tel que
$|\varphi(s)-\varphi(t)|\le\eta$ et $|\varphi'(s)-\varphi'(t)|\le\eta$
pour $s^{-1}t\in V$. Soit alors $g\in\mathcal K_+^*(V)$ ; pour tout $s\in G$,
on a $\varphi.\gamma(s)g\le(\varphi(s)+\eta).\gamma(s)g$
(en effet, c’est évident aux points où $\gamma(s)g$ s’annule, donc hors de $sV$ ;
et, dans $sV$, $\varphi\le\varphi(s)+\eta$) ; de même,
$\varphi'.\gamma(s)g\le(\varphi'(s)+\eta).\gamma(s)g$.
Ceci posé, soient $c_1,\ldots,c_n$ des nombres $\ge0$ et
$s_1,\ldots,s_n$ des éléments de $G$ tels que
$F\le\sum_{i=1}^n c_i\gamma(s_i)g$ ; on a
\[
f=\varphi F\le\sum_{i=1}^n c_i\varphi.\gamma(s_i)g
\le\sum_{i=1}^n c_i(\varphi(s_i)+\eta).\gamma(s_i)g
\]
et de même pour $f'$ ; par suite
\[
(f:g)+(f':g)\le\sum_{i=1}^n c_i(\varphi(s_i)+\varphi'(s_i)+2\eta)
\le(1+2\eta)\sum_{i=1}^n c_i
\]
puisque $\varphi+\varphi'\le1$. En appliquant la définition de $F$,
puis (ii), (iii), et (v), on en conclut
\[
\begin{aligned}
(f:g)+(f':g)&\le(1+2\eta)(F:g)
\le(1+2\eta)\bigl[(f+f':g)+\tfrac12\varepsilon(h:g)\bigr]\le{}\\
&\le(f+f':g)+\tfrac12\varepsilon(h:g)
+2\eta(f+f':h)(h:g)+\varepsilon\eta(h:g)
\end{aligned}
\]
et, si on a choisi $\eta$ tel que
$\eta[2(f+f':h)+\varepsilon]\le\frac12\varepsilon$, on obtient (vii).

Quand $V$ parcourt l’ensemble des voisinages compacts de $e$, les
$\mathcal K_+^*(V)$ forment une base de filtre $\mathcal B$ sur $\mathcal K_+^*$.
Soit $\mathcal F$ un ultrafiltre sur $\mathcal K_+^*$ plus fin que $\mathcal B$.
